\section{Problem Formulation}
\label{sec:problem-formulation}

\subsection{Domain and three linked questions}

Both optimization questions use the prescribed Kaiwu workflow.  We
distinguish three execution paths throughout the paper.  \texttt{SDK--CPU}
refers to local CPU sampling with Kaiwu SDK~1.4.3.  \texttt{Mapping audit}
refers to Kaiwu~1.4.4 conversion and field-compatibility checks; no solve
licence was available for that version.  \texttt{SPQC-550} refers only to
platform tasks with a recorded task ID, submitted-matrix hash, and archived
return.  We do not transfer performance evidence between these paths.

\subsection{Violation score and optimization objectives}

For a signed deviation $d$ in eV, define
\begin{equation}
 \phi(d)=\left(\frac{\max(0,d)}{0.5}\right)^2 .
 \label{eq:formulation-hinge}
\end{equation}
With formation energy $e_{\rm form}$, band gap $E_g$, and band edges
$E_{\rm VBM}$ and $E_{\rm CBM}$, the official score is
\begin{align}
 V(z)={}&\phi(e_{\rm form}-0.21)+\phi(1.4-E_g)+\phi(E_g-3.1)\notag\\
       &+\phi(5.73-E_{\rm VBM})+\phi(E_{\rm CBM}-4.50).
 \label{eq:formulation-V}
\end{align}
The squared-hinge score is zero at equality with a threshold, whereas
strict feasibility requires
$e_{\rm form}<0.21$, $1.4<E_g<3.1$, $E_{\rm VBM}>5.73$, and
$E_{\rm CBM}<4.50$.  We therefore report the score and strict feasibility
as separate quantities; a low predicted score does not certify that a
candidate satisfies the strict conditions.

For Q1, model selection uses files~02 and~03.  After the specification
is frozen, the release artifact \texttt{q1\_final} is fitted on file~02
alone and applied to the unlabeled file~04.  We assess the selected model
using MAE, RMSE, Spearman correlation, low-tail error, and Top-$K$
retrieval.  For Q2/Q3, the frozen architecture is refitted on files~02
and~03 to score the candidate pools.  We denote this separate score by
$\widehat V_{23}$ and define the two optimization objectives as
\begin{align}
 z_2^*&=\arg\min_{A,B\in\mathcal E}\widehat V_{23}(A,B,\mathrm{NO}_2),
 \label{eq:q2-objective}\\
 z_3^*&=\arg\min_{A,B\in\mathcal E,\,X\in\mathcal X_3}
 \widehat V_{23}(A,B,X).
 \label{eq:q3-objective}
\end{align}
They are frozen-proxy objectives, not hidden-label measurements.

\subsection{Binary interface and groupwise QUBO}

Index $\mathcal E$ by $i,j=1,\ldots,52$ and introduce one-hot variables
$a_i,b_j\in\{0,1\}$ satisfying
\begin{equation}
 \sum_i a_i=1,\qquad \sum_j b_j=1.
 \label{eq:one-hot-constraints}
\end{equation}
For fixed $X$, let $c_{ij}^{(X)}$ be the traceable nonnegative integer cost
derived from candidate $(i,j,X)$.  The adopted 104-variable mapping is
\begin{align}
 E_X(\bm a,\bm b)={}&\sum_{i,j}c_{ij}^{(X)}a_i b_j
 +\lambda_A\left(\sum_i a_i-1\right)^2
 +\lambda_B\left(\sum_j b_j-1\right)^2 .
 \label{eq:group-qubo}
\end{align}
Expanding the constraints gives a QUBO
\cite{Glover2019QUBOTutorial,Lucas2014IsingNP}.  Q2 uses only the NO2
matrix.  Q3 solves three such matrices and compares decoded winners on the
common original scale,
\begin{equation}
 X^*=\arg\min_{X\in\mathcal X_3}
 \widehat V_{23}(A_X^*,B_X^*,X).
 \label{eq:q3-group-comparison}
\end{equation}
The proxy enumeration covers all 8,112 Q3 candidates and is distinct from the
earlier 208-variable joint control.  Chapters~\ref{sec:q2-qubo}--\ref{sec:hardware-precision}
audit the cost transform, offset, coefficient bounds, and decoding.

\subsection{Assumptions and evidence boundary}

The notation used below is as follows: $V$ denotes the observed
competition score, $\widehat V$ the Q1 surrogate prediction,
$\widehat V_{23}$ the separately refitted Q2/Q3 prediction, and $E_X$ the
encoded energy for a fixed anion group.

Our working scope follows the organizer's candidate grid and score
definition.  The A and B sites are ordered, $A=B$ is allowed, and only
composition fields and deterministic composition features are used;
identifiers, row order, pool membership, and physical labels are excluded
from model inputs.  The supplied six-decimal values are used as given and
do not relax the strict feasibility inequalities.  Because the Q2/Q3
pools are unlabeled, their optimization objective is the frozen proxy
score, not the unobserved physical score.  Only assignments satisfying
both one-hot constraints are treated as legal candidates.  These choices
define the computational scope but do not remove surrogate error.

We keep surrogate prediction, QUBO equivalence, coefficient adaptation,
solver sampling, and physical validation as separate sources of
uncertainty.  Exhaustive proxy enumeration identifies the minimum of the
frozen score table; encoded-state enumeration checks whether the QUBO
represents that known minimum.  Neither step establishes material
feasibility.  An SPQC-550 return is likewise a sample from the encoded
proxy objective, not a measurement of $V$ or a certification of the
selected material.
